\documentclass[11pt,a4paper,twocolumn]{article}

% --- Packages ---
\usepackage[utf8]{inputenc}
\usepackage[T1]{fontenc}
\usepackage[english]{babel}
\usepackage{amsmath,amssymb}
\usepackage{booktabs}
\usepackage{graphicx}
\usepackage{xcolor}
\usepackage{hyperref}
\usepackage{cleveref}
\usepackage{float}
\usepackage{enumitem}
\usepackage{caption}
\usepackage{subcaption}
\usepackage{geometry}
\usepackage{authblk}
\usepackage{abstract}
\usepackage{tikz}
\usepackage{pgfplots}
\pgfplotsset{compat=1.18}
\usepackage{multirow}
\usepackage{tabularx}
\usepackage{colortbl}
\usepackage{csquotes}
\usepackage{epigraph}

\geometry{margin=2cm,top=2.5cm,bottom=2.5cm}

\hypersetup{
    colorlinks=true,
    linkcolor=blue!60!black,
    citecolor=blue!60!black,
    urlcolor=blue!60!black
}

% --- Custom colors ---
\definecolor{cinit}{HTML}{378ADD}
\definecolor{cclose}{HTML}{E24B4A}
\definecolor{clink}{HTML}{EF9F27}
\definecolor{cact}{HTML}{639922}
\definecolor{cmode}{HTML}{7F77DD}
\definecolor{ctime}{HTML}{1D9E75}
\definecolor{cref}{HTML}{888780}
\definecolor{highlight}{HTML}{FFF8E1}
\definecolor{deepblue}{HTML}{0C447C}

% --- Version ---
\newcommand{\PaperVersion}{v0.30}

% --- Title ---
\title{\textbf{Structured Domain-Specific Notation\\in the Voynich Manuscript:}\\[4pt]
\large A Multi-Layer Recording Architecture for Observational\\
and Procedural Knowledge --- Draft \PaperVersion}

\author[1]{Roble Mumin\thanks{Corresponding author. \url{https://roblemumin.com}}}
\affil[1]{Independent AI Researcher and AI Consultant}

\date{March 2026 --- Draft \PaperVersion}

\begin{document}

\maketitle

\epigraph{\textit{A notation system is not the knowledge it encodes.\\
It is the architecture that makes that knowledge transmissible.}}{--- Mumin, 2026c}

% ============================================================
\begin{abstract}
\noindent Two companion papers (Mumin, 2026a; 2026b) established that the Voynich Manuscript (Beinecke MS 408, c.\ 1404--1438) encodes a six-role positional functional system organized into complete packet structures, with quantifiably distinct domain profiles across six section types. The present paper advances an interpretive thesis grounded in those structural findings: we argue that the Voynich Manuscript is the earliest surviving example of a \textit{highly structured, multi-layer domain-specific notation system} designed for concise and repeatable recording of observational and procedural knowledge in late-medieval medicine, astrology, and pharmacy. This thesis does not require the system to encode natural language. We demonstrate that the Voynich text exhibits four defining properties of domain-specific notation: (1)~rigid positional syntax enforced across all sections, (2)~a finite-state packet grammar with 61.3\% paragraph-level conformance, (3)~a three-layer abstraction architecture separating structural frame, inner-packet function words, and lexical entities, and (4)~a folio-anchor labeling pattern consistent with per-entity referencing rather than narrative discourse. We compare this architecture with contemporaneous notation traditions in liturgical music, astronomical tables, and pharmaceutical recipe notation, showing that the Voynich system surpasses all of them in structural formality and layer separation. We propose that the system was produced by a single specialist or small specialist community operating outside mainstream clerical culture---and we briefly discuss, as a speculative epilogue clearly labeled as such, the possibility that the specialist was non-human. We argue that the notation-history framing resolves several long-standing puzzles about the manuscript without requiring any assumption about its content or decipherability. This paper does not determine what the system encodes.

\medskip
\noindent\textbf{Keywords:} Voynich Manuscript, domain-specific notation, recording systems, functional grammar, packet architecture, notation history, medieval manuscript studies, structural analysis
\end{abstract}

% ============================================================
\section{Introduction}
\label{sec:intro}

\subsection{The Question of Interpretive Frame}

After six centuries of failed decipherment attempts, the Voynich Manuscript (Beinecke MS 408) remains undeciphered. The dominant interpretive frame has been linguistic: most investigators have assumed the manuscript encodes a natural language, either directly or through a cipher, and have applied the tools of linguistics and cryptanalysis accordingly. When those tools fail, the manuscript is labeled anomalous, random, or a hoax.

The companion papers in this series (Mumin, 2026a; 2026b) took a different approach. Rather than assuming linguistic content, they asked: what kind of structural system does the Voynich text actually exhibit? The results were definite. The text encodes a six-role positional functional system with significant directional asymmetry, organized into complete packet structures, with coherent domain-specific distributions across section types. These are properties of a \textit{recording system}---but not necessarily a linguistic one.

The present paper asks a further question: \textit{what kind of recording system?} We propose that the most accurate and parsimonious answer is a \textit{domain-specific notation system} (DSN): a structured symbolic code designed not for general communication but for the efficient, repeatable capture of domain-specific observational and procedural knowledge.

This framing resolves several puzzles that the linguistic frame cannot: why the text has no narrative structure, no argument relations, no embedded clauses; why each domain section has a distinct functional profile; why the system uses a three-layer abstraction rather than a flat symbol alphabet; and why it is internally consistent over more than 37,000 tokens without a single decoded word.

\subsection{Scope and Limitations}

This paper makes structural and comparative claims. It does not:

\begin{itemize}[leftmargin=*,itemsep=2pt]
    \item Identify the language underlying the notation, if any
    \item Propose specific token readings or decipherment candidates
    \item Claim to determine what specific knowledge the system encodes
    \item Resolve the question of authorship beyond the structural evidence
\end{itemize}

\noindent The claims it does make are conditional on the structural findings in Mumin (2026a) and Mumin (2026b), which are themselves subject to the limitations stated in those papers: the six-role classification is one of several competitive models, and the packet grammar describes a statistical tendency rather than an algorithmic constraint.

\subsection{Organization}

Section~\ref{sec:structural} summarizes the structural foundations from Papers 1 and 2. Section~\ref{sec:dsn} defines domain-specific notation and its distinguishing properties. Section~\ref{sec:vm_as_dsn} demonstrates that the Voynich system satisfies all four defining properties. Section~\ref{sec:comparison} compares the Voynich architecture with contemporaneous notation traditions. Section~\ref{sec:specialist} evaluates the specialist-producer hypothesis. Section~\ref{sec:epilogue} presents the extraterrestrial hypothesis as an explicitly speculative epilogue. Section~\ref{sec:future} outlines a research program, and Section~\ref{sec:conclusion} concludes.

% ============================================================
\section{Structural Foundations}
\label{sec:structural}

\subsection{Six-Role Positional System (Paper 1)}

Mumin (2026a) established that Voynich EVA glyph clusters occupy statistically distinct positional roles within word-level tokens. The six-role taxonomy---derived from position-weighted frequency profiling and validated by role-inversion falsification---characterizes the following functional categories:

\begin{table}[H]
\centering
\caption{Six-role functional taxonomy (Paper 1)}
\label{tab:roles}
\small
\begin{tabular}{@{}clp{3.5cm}@{}}
\toprule
\textbf{Role} & \textbf{Label} & \textbf{Defining Feature} \\
\midrule
R1 & \textsc{init} & Packet-initiating; \texttt{qok-} prefix family \\
R2 & \textsc{close} & Packet-closing; \texttt{ch/sh/lch} stems \\
R3 & \textsc{link} & Transitional; inter-packet connectors \\
R4 & \textsc{content} & High-frequency domain-specific tokens \\
R5 & \textsc{temporal} & Low-frequency, section-specific \\
R6 & \textsc{ref} & Reference markers; \texttt{ol/al} family \\
\bottomrule
\end{tabular}
\end{table}

\noindent The critical validation result is the R2$\to$R1 transition: this boundary (packet close followed by packet open) is the strongest positional asymmetry in the corpus ($z = +9.71$), surviving both role inversion ($z < 1.0$ under shuffled assignments) and section stratification. The asymmetry demonstrates that the positional roles are not artifacts of the clustering procedure but structural features of the notation.

\subsection{Packet Grammar and Domain Coherence (Paper 2)}

Mumin (2026b) extended the six-role model to the \textit{packet level}: sequences of functional roles from R1 (init) through payload to R2 (close), with R6 (ref) tokens acting as frame markers. Key quantitative results:

\begin{itemize}[leftmargin=*,itemsep=2pt]
    \item 941 complete packets detected (R1$\to$payload$\to$R2) in the sample corpus
    \item 61.3\% of paragraphs conform to expected FSA structure
    \item 52.1\% of corpus tokens fall within the R1/R2/R6 structural frame
    \item Section-level domain profiles are statistically distinct across all six section types (Herbal, Astronomical, Balneological, Pharmaceutical, Recipe, Stars)
\end{itemize}

\noindent The domain coherence finding is particularly relevant for the DSN thesis: the system's functional profile shifts systematically by section, which is a defining feature of domain-specific rather than general-purpose notation.

\subsection{Three-Layer Architecture (PILOT3)}

A further structural result, established in PILOT3 (this series), identifies three separable layers within the packet architecture:

\begin{enumerate}[leftmargin=*,itemsep=2pt]
    \item \textbf{Structural frame}: \texttt{ol} (R6-classified), elevated inside complete packets at 1.41$\times$ the rate expected by chance
    \item \textbf{Inner-packet function words}: \texttt{qol}, elevated at 1.39$\times$ inside packets; paragraph co-occurrence with \texttt{ol}: $\mathrm{OR} = 64.6$, $p < 0.0001$; not positionally biased within payloads (mean position 0.558, $p = 0.23$)
    \item \textbf{Lexical entities}: 44 Rosetta-candidate tokens with section-specific distributions, folio-anchoring patterns, and co-occurrence constraints (e.g., \texttt{sal}+\texttt{oly} co-occurrences $= 0$)
\end{enumerate}

\noindent The near-identity of \texttt{qol} (1.39$\times$) and \texttt{ol} (1.41$\times$) enrichment inside packets is the key finding: an inner-packet function word behaves like the structural frame marker, suggesting a grammatical parallel between layers. Note that \texttt{kol} (0.42$\times$) and \texttt{okol} (0.73$\times$)---despite their surface similarity to \texttt{qol}---are \textit{depleted} inside packets, confirming that the three-layer separation is not a simple prefix-variation artifact.

% ============================================================
\section{Domain-Specific Notation: Definition and Properties}
\label{sec:dsn}

\subsection{What Is Domain-Specific Notation?}

A \textit{domain-specific notation system} (DSN) is a symbolic code designed for the efficient, repeatable representation of knowledge within a constrained domain---as opposed to a \textit{general-purpose language} capable of expressing arbitrary propositions. The concept is well-established in modern computing (domain-specific languages, DSLs), but the underlying design principle is ancient: musical notation records pitch and rhythm efficiently but cannot express propositions; astronomical tables encode positional data in compact form but carry no narrative; pharmaceutical recipe notation captures ingredient, quantity, and procedure without argumentation.

Three properties distinguish DSNs from natural languages and general-purpose ciphers:

\begin{enumerate}[leftmargin=*,itemsep=2pt]
    \item \textbf{Positional rigidity}: Symbol roles are determined by position, not by inflection or agreement. The grammar is a structure constraint, not a generative syntax.
    \item \textbf{Finite state}: The system can be described by a finite-state automaton over a small alphabet of functional categories. It does not require recursion, embedding, or argument structure.
    \item \textbf{Domain coupling}: The functional profile of the notation shifts with the domain being recorded. A general-purpose language has the same grammar across all topics; a DSN has domain-specific layers.
\end{enumerate}

\noindent A fourth property---\textbf{entity labeling rather than narrative}---is common but not universal: many DSNs reference entities (stars, plants, procedures) rather than making claims about them.

\subsection{Why These Properties Matter for the Voynich Debate}

The linguistic frame predicts that the Voynich system should exhibit argument structure, word-order variability consistent with natural language parameters, and statistical properties (Zipf slope, type-token ratio) consistent with language. It predicts none of these are observed without invoking complex cipher transformations.

The DSN frame predicts exactly what is observed: rigid positional roles, finite-state packet grammar, domain-specific functional profiles, and statistical properties atypical of natural language (Zipf exponent 0.555, hapax ratio 0.780---both consistent with a constrained technical vocabulary). The DSN frame requires no additional assumptions to account for these features.

% ============================================================
\section{The Voynich System as Domain-Specific Notation}
\label{sec:vm_as_dsn}

We now demonstrate that the Voynich text satisfies each of the four defining properties of DSN.

\subsection{Property 1: Positional Rigidity}

The six-role taxonomy established in Paper 1 is fundamentally a positional taxonomy: R1 tokens appear at packet starts, R2 tokens at packet ends, R6 tokens as frame markers, and so on. This positional assignment is statistically robust across all section types and survives the role-inversion falsification test.

The rigidity is stronger than any known natural language word-order constraint. In natural languages with strict verb-second (V2) order---the tightest attested positional constraint---the relevant unit is a grammatical constituent (verb phrase), and the constraint tolerates a wide range of lexical fillers. In the Voynich system, the role assignments are cluster-specific: individual token types occupy their roles with high fidelity (R1: $\kappa = 0.72$; R2: $\kappa = 0.71$; R6: $\kappa = 0.68$ in the full-corpus run).

This level of positional rigidity is characteristic of notation, not language.

\subsection{Property 2: Finite-State Packet Grammar}

The SPBCEH packet grammar is formally a finite-state automaton over the six-role alphabet. The dominant transition structure is:

\begin{center}
\texttt{R1 $\to$ (R4|R6)* $\to$ R2}
\end{center}

\noindent with R3 (link) tokens appearing between packets and R5 (temporal) tokens appearing rarely within payloads. This structure is finite: it requires no recursion, no embedded clauses, no memory of unbounded depth.

The 61.3\% paragraph-level FSA conformance rate confirms that the majority of the corpus follows this template. The 38.7\% non-conformance includes packet fragments (truncation), apparent scribal variations, and section-specific deviations---all consistent with notation use rather than linguistic generation.

No natural language has a finite-state grammar. Human languages require at minimum context-free power (center embedding, cross-serial dependencies). The Voynich grammar does not. This is a decisive structural difference between the Voynich system and any natural language, coded or uncoded.

\subsection{Property 3: Domain-Specific Functional Profiles}

The six section types of the manuscript exhibit quantifiably distinct R4 (\textsc{content}) token distributions:

\begin{table}[H]
\centering
\caption{Domain-specific functional divergence by section}
\label{tab:domain}
\small
\begin{tabular}{@{}lcc@{}}
\toprule
\textbf{Section} & \textbf{R4 Entropy} & \textbf{Top R4 Token} \\
\midrule
Herbal (H) & High & \texttt{daiin} \\
Astronomical (A) & Medium & \texttt{ol} \\
Balneological (B) & Low & \texttt{daiin} (dense) \\
Pharmaceutical (P) & High & \texttt{qokeey} \\
Recipe (R) & Medium & \texttt{cheol} \\
Stars (S) & Low & \texttt{-ain} family \\
\bottomrule
\end{tabular}
\end{table}

\noindent The domain-specific divergence is not a statistical artifact: the section-classification experiment (Paper 2) correctly assigns folios to their section type from role profiles alone at rates significantly above chance ($p < 0.001$, 6.9 pp gain over baseline). A general-purpose cipher would show no domain-specific functional differentiation.

The Stars section exhibits an additional domain-coupling feature: dominant \texttt{-ain} family tokens vary systematically across folios (PILOT2). While the strong entity-labeling hypothesis (dominant \texttt{-ain} token encodes illustration subject) is not supported---same-subject recto/verso consistency is only 2/7 (29\%), and Arabic constellation alignment scores range from 0.078 to 0.333---the descriptive pattern is real: each Stars folio has a distinct dominant \texttt{-ain} cluster, and this cluster differs across folios. This is consistent with sub-entity labeling (e.g., individual star groups or nymph figures within a zodiac circle) rather than simple zodiac-sign encoding.

\subsection{Property 4: Entity Labeling Rather Than Narrative}

The folio-anchor \texttt{-ain} pattern in the Stars section, the 44 Rosetta-candidate entity tokens with section-specific distributions, and the structural separation of the three layers (frame / function words / lexical entities) all point toward entity-referencing rather than narrative discourse.

In the three-layer architecture, layer 3 (lexical entities) corresponds to what a DSN calls a \textit{symbol vocabulary}: a finite set of referential tokens that name the entities being recorded (plants, stars, procedures, anatomical regions) without asserting propositions about them. The structural separation of entity tokens from functional frame and inner-function tokens is the architectural signature of this design: the system \textit{can} enumerate entities without predicating anything of them.

The zero co-occurrence of \texttt{sal} and \texttt{oly}---two candidate entity labels with high section-specificity---is consistent with non-overlapping entity domains: \texttt{sal} appears in contexts where \texttt{oly} does not, suggesting these tokens label distinct entity classes that the system treats as categorically exclusive.

\textbf{Caveat}: The entity-labeling interpretation of layer 3 is a structural hypothesis, not a confirmed reading. Confirming it would require successful alignment between candidate entity tokens and known domain objects (plant names, star names, ingredient names) across multiple independent sources.

% ============================================================
\section{Comparison with Medieval Notation Traditions}
\label{sec:comparison}

\subsection{Contemporaneous Notation Systems}

The 15th century had a rich landscape of specialized notation systems. We compare the Voynich architecture with four:

\begin{table}[H]
\centering
\caption{Structural comparison with contemporaneous notation traditions}
\label{tab:comparison}
\small
\begin{tabular}{@{}p{2.0cm}cccc@{}}
\toprule
\textbf{System} & \rotatebox{65}{\textbf{Pos.\ rigid}} & \rotatebox{65}{\textbf{Finite-state}} & \rotatebox{65}{\textbf{Domain-coupled}} & \rotatebox{65}{\textbf{Entity-label}} \\
\midrule
Mensural music & \checkmark & \checkmark & --- & --- \\
Astron.\ tables & \checkmark & $\sim$ & \checkmark & \checkmark \\
Recipe notation & --- & $\sim$ & \checkmark & $\sim$ \\
Herbal notation & --- & --- & \checkmark & \checkmark \\
\textbf{Voynich (SPBCEH)} & \checkmark & \checkmark & \checkmark & \checkmark \\
\bottomrule
\end{tabular}
\end{table}

\subsection{Mensural Musical Notation}

By the 15th century, mensural notation (Ars Nova through the Burgundian School) had developed a sophisticated positional system: noteheads encode pitch by staff position, note shapes encode duration, and ligatures encode groupings. The system is positionally rigid and finite-state (note sequences are constrained by mensural rules), but it is domain-limited to pitch-rhythm recording and does not couple its functional profile to the subject domain.

The Voynich system surpasses mensural notation in one key respect: three-layer abstraction. Mensural notation has no equivalent to the Voynich's structural-frame / inner-function-word / lexical-entity separation.

\subsection{Astronomical Tables}

Medieval astronomical tables (such as the \textit{Alfonsine Tables}, c.\ 1270, still in heavy use in the 15th century) are among the most formally structured medieval notations. They encode planetary positions, eclipse predictions, and calendrical data in tabular form with positional rigidity: column order encodes variable type, row structure encodes the enumerated entity, and numerical entries are the observational values.

Astronomical tables are entity-referencing rather than narrative, and they are domain-coupled. However, they are not finite-state in the Voynich sense---they do not have a grammar; they are purely relational. The Voynich's packet grammar is a structural sophistication that astronomical tables lack.

\subsection{Pharmaceutical Recipe Notation}

Medieval recipe notation (\textit{receptaria}, \textit{antidotaria}) encodes ingredient, quantity, preparation method, and indication in roughly consistent order. However, the order is flexible (recipe notation tolerates inversion and omission of components), and the vocabulary is natural language---Latin or vernacular---with full propositional structure in the indication field. Recipe notation is domain-coupled but not positionally rigid and not finite-state.

\subsection{Herbal and Botanical Notation}

Medieval herbals (such as the \textit{Tractatus de Herbis} tradition) pair botanical illustrations with verbal descriptions in natural language. The illustration-text pairing is structurally fixed, and the illustrations are entity-referencing, but the text is propositional natural language. The Voynich system's text layer lacks all propositional features.

\subsection{What the Voynich System Adds}

The Voynich system is the only 15th-century notation that simultaneously exhibits all four defining properties of DSN. It is more formally structured than any contemporaneous specialty notation:

\begin{itemize}[leftmargin=*,itemsep=2pt]
    \item More positionally rigid than recipe notation
    \item More grammatically constrained than astronomical tables
    \item More domain-differentiated than musical notation
    \item More architecturally layered than any of them
\end{itemize}

\noindent If the structural analysis is correct, the Voynich Manuscript represents an advance in formal notation design that was not matched by any other known medieval system---and that anticipates properties of modern formal languages by several centuries.

% ============================================================
\section{The Specialist-Producer Hypothesis}
\label{sec:specialist}

\subsection{Who Could Have Designed This System?}

The structural formality of the Voynich notation implies a producer with:

\begin{enumerate}[leftmargin=*,itemsep=2pt]
    \item \textbf{Domain expertise across multiple fields}: The system covers at least six distinct knowledge domains (herbal medicine, astronomy, balneology, pharmacy, recipes/procedures, stellar observation) with domain-specific functional profiles for each. This breadth is characteristic of a generalist specialist---a physician-astrologer, a court natural philosopher, or a monastic polymath---not a single-domain practitioner.
    \item \textbf{Formal notation design experience}: The three-layer architecture and finite-state packet grammar reflect deliberate structural design rather than ad-hoc notation. The producer understood the principle of separating structural frame from content from reference---a design insight not explicitly articulated in any 15th-century notation treatise but demonstrably implemented here.
    \item \textbf{Motivation for non-linguistic encoding}: A producer working in mainstream clerical culture would default to Latin or a vernacular language. Non-linguistic notation implies either (a) the need for compactness beyond what natural language allows, (b) the desire for concealment from non-initiated readers, (c) a multi-lingual community that lacked a shared natural language for the domain, or (d) a producer whose primary cognitive medium was not a natural language.
    \item \textbf{Extended solitary or small-group production}: The internal consistency of the notation across 37,000+ tokens implies sustained, coherent production---whether by a single author or a small group following a shared notation convention. No evidence of multiple competing notation variants has been identified.
\end{enumerate}

\subsection{The Single-Author Efficiency Argument}

The three-layer architecture is most parsimonious as a single-author design. In a collaborative or transmitted notation system, we would expect competing variants, layer-boundary ambiguities, and gradual evolution across the manuscript. The consistency of \texttt{qol} behavior (1.39$\times$ enrichment), \texttt{ol} behavior (1.41$\times$), and the \texttt{sal}/\texttt{oly} exclusion pattern across independent sub-corpora suggests a unified design rather than a collaborative accretion.

The closest historical parallel is the design of formal shorthand systems by individual practitioners: Tironian notes (attributed to Marcus Tullius Tiro), Petrus Balesdens's 14th-century shorthand variants, and later systems by Timothy Bright (1588) and John Willis (1602) were all individual inventions with the kind of systematic internal coherence observed here.

\subsection{Why Concealment Alone Is Insufficient}

The most common specialist-producer hypothesis is that the Voynich system is a cipher designed to conceal its content from unauthorized readers. This is consistent with the evidence but insufficient as a complete explanation.

A cipher over a natural language would retain the statistical properties of that natural language under most cipher transformations. The Voynich's Zipf exponent (0.555, vs.\ $\sim$1.0 for natural languages) and its hapax ratio (0.780) are not consistent with any known cipher of a natural language, unless the cipher involves dramatic vocabulary compression---which would require a second notation layer on top of the natural language.

The DSN hypothesis is more parsimonious: the anomalous statistics arise not from cipher transformation but from the original design of the notation, which was built for efficiency (short tokens, constrained vocabulary) rather than expressiveness.

% ============================================================
\section{Speculative Epilogue: The Non-Human Producer Hypothesis}
\label{sec:epilogue}

\textbf{This section is explicitly speculative.} The following argument is presented as a candidate hypothesis that follows from the DSN framing under one set of additional assumptions. It is not supported by direct evidence and should not be read as a claim of the same evidential status as the structural findings in Sections~\ref{sec:structural}--\ref{sec:specialist}.

\subsection{From Unusual Specialist to Non-Human Producer}

The specialist-producer hypothesis (Section~\ref{sec:specialist}) requires four conditions: multi-domain expertise, formal notation design experience, motivation for non-linguistic encoding, and extended solitary production. For a human producer, condition (4) is the most demanding: 240+ pages of internally consistent notation with no natural-language features, no identifiable cultural referents in the illustrations, and no conventional scribal colophons or marginalia. Human specialists who produce extended private notation systems typically leave some natural-language trace---a title, a label, a marginal correction in the vernacular. The Voynich Manuscript has none.

One explanation for this absence is a producer for whom natural language was not the default cognitive medium---whose primary representational system was the notation itself rather than a natural language that the notation encodes or conceals. Such a producer would not default to natural-language traces because natural language was not the default.

This property is most easily explained by a producer who is \textit{not a language-using species}---that is, a non-human intelligence. Under the DSN framing, the Voynich system is that producer's domain-specific notation, designed from scratch to capture observational knowledge in a medium (vellum, iron gall ink) that the producer acquired locally.

\subsection{What the DSN Framing Adds to the ET Hypothesis}

The version of the non-human authorship hypothesis advanced in the companion paper to this series (Mumin, 2026c, earlier draft) rested primarily on negative evidence: the manuscript cannot be explained by human cognitive architectures, therefore it may be non-human. The DSN framing adds a positive contribution: the system is not merely anomalous; it is \textit{sophisticated in a specific way} that is characteristic of deliberate formal notation design.

A non-human intelligence designing a recording system from scratch, without cultural templates for notation, might converge on exactly this architecture: positionally rigid (because positional encoding requires no shared cultural convention), finite-state (because finite-state systems are the minimal formal structure needed for structured recording), domain-coupled (because different knowledge domains require different functional vocabularies), and entity-referencing (because the goal is to record entities and their properties, not to make propositional arguments about them).

In this reading, the Voynich system is not a degraded or encrypted version of a human notation. It is an \textit{independently designed} notation that converges on formal properties that human notation systems took centuries to develop---because those properties are, in some sense, optimal for structured knowledge recording.

\subsection{Status of This Hypothesis}

The non-human producer hypothesis is:

\begin{itemize}[leftmargin=*,itemsep=2pt]
    \item \textbf{Consistent} with the structural evidence from Papers 1--3
    \item \textbf{Not supported} by direct physical or palaeographic evidence
    \item \textbf{Falsifiable} at multiple levels (perceptual filtering prediction, material anomaly prediction, cross-artifact search)
    \item \textbf{Not the primary thesis} of this paper or this series
\end{itemize}

\noindent We include it here because intellectual honesty requires acknowledging that the structural evidence, taken to its logical limit, does not exclude it---and because the DSN framing provides a more precise formulation of the hypothesis than the ``cognitive state log'' framing in earlier drafts. Whether it merits serious investigation is a judgment we leave to readers.

% ============================================================
\section{Research Program}
\label{sec:future}

The DSN framing generates specific, testable research directions distinct from those implied by the linguistic frame.

\subsection{Notation Grammar Formalization}

\begin{itemize}[leftmargin=*,itemsep=2pt]
    \item Formalize the SPBCEH packet grammar as a complete finite-state automaton over all 37,045 corpus tokens
    \item Test grammar conformance rates across section types and folio boundaries
    \item Identify systematic deviations that may represent notation evolution or scribal variants
\end{itemize}

\subsection{Layer 3 Entity Alignment}

\begin{itemize}[leftmargin=*,itemsep=2pt]
    \item Extend the 44 Rosetta-candidate entity vocabulary
    \item Test alignment between candidate entity tokens and domain-specific referent lists: plant names (Latin, Arabic, vernacular), star names, pharmaceutical ingredients, anatomical terms
    \item Use co-occurrence constraints (e.g., \texttt{sal}/\texttt{oly} exclusion) to infer entity-category structure
    \item Evaluate multiple languages simultaneously (Latin, Hebrew, Arabic, vernacular Italian/German) as potential layer 3 alignment targets
\end{itemize}

\subsection{Stars Section Sub-Entity Analysis}

\begin{itemize}[leftmargin=*,itemsep=2pt]
    \item The PILOT2 finding (folio-anchor pattern real; strong entity-labeling hypothesis not supported; recto/verso consistency 2/7) suggests that \texttt{-ain} tokens label sub-folio entities (individual star groups or nymph figures) rather than zodiac signs
    \item Test this by mapping dominant \texttt{-ain} tokens against illustration sub-elements within each folio
    \item If confirmed, this would establish the finest granularity of entity labeling in the system
\end{itemize}

\subsection{Notation History Contextualization}

\begin{itemize}[leftmargin=*,itemsep=2pt]
    \item Survey 14th--15th century shorthand, recipe, astronomical, and medical notation systems for structural features that anticipate the Voynich architecture
    \item Identify candidate inventor communities (court physicians, monastic scholars, alchemists, astrologers) whose known notation practices most closely approach the Voynich design
    \item Search for traces of the Voynich notation system in undescribed manuscripts from the same period
\end{itemize}

\subsection{Structural Comparison with Other Undeciphered Scripts}

\begin{itemize}[leftmargin=*,itemsep=2pt]
    \item Apply the SPBCEH role-classification framework to other undeciphered texts (Linear A, Proto-Elamite, Indus script) to determine whether any exhibit comparable positional rigidity, finite-state grammar, and three-layer architecture
    \item If any do, this suggests the Voynich architecture is not unique but represents a convergent design solution for structured domain-specific recording
\end{itemize}

% ============================================================
\section{Conclusion}
\label{sec:conclusion}

We have argued that the Voynich Manuscript is the earliest surviving example of a structured, multi-layer domain-specific notation system: a symbolic code designed for the efficient, repeatable recording of observational and procedural knowledge in medicine, astronomy, and pharmacy, operating without natural-language propositional structure.

This thesis rests on four structural demonstrations:

\begin{enumerate}[leftmargin=*,itemsep=2pt]
    \item The Voynich text exhibits \textit{positional rigidity} surpassing any known natural-language word-order constraint ($\kappa = 0.68$--$0.72$ per role)
    \item The packet grammar is \textit{finite-state} (61.3\% paragraph conformance), placing it outside the formal power required by any natural language
    \item The functional profile is \textit{domain-coupled}, shifting coherently across six section types in a pattern consistent with domain-specific knowledge recording
    \item The three-layer architecture (\texttt{ol} frame / \texttt{qol} inner-function / 44 Rosetta entities) is an \textit{entity-referencing} rather than propositional design
\end{enumerate}

\noindent Compared with contemporaneous notation traditions---mensural music, astronomical tables, pharmaceutical recipes, botanical herbals---the Voynich system is the most formally complete: the only one exhibiting all four DSN properties simultaneously.

The specialist-producer hypothesis (a single domain-expert author operating outside mainstream clerical culture) is the most parsimonious explanation for the system's internal consistency, multi-domain scope, and absence of natural-language residue. We have presented, in a clearly labeled speculative epilogue (Section~\ref{sec:epilogue}), the stronger hypothesis that the producer was non-human; we do not claim this is established.

\medskip

What this paper does not establish:

\begin{itemize}[leftmargin=*,itemsep=2pt]
    \item What specific knowledge the notation encodes
    \item What language, if any, underlies the layer 3 entity vocabulary
    \item Whether the system is in principle decipherable
    \item Who produced it
\end{itemize}

\noindent The notation-history framing earns its place not by answering these questions but by explaining, without auxiliary assumptions, why six centuries of linguistic and cryptanalytic approaches have not. The Voynich Manuscript may be undeciphered not because it is too complex for human analysis, but because it is not the kind of artifact that linguistic analysis is designed to decode.

\medskip
\noindent\textbf{What this paper does establish:} The Voynich Manuscript's text layer is a structured, internally consistent, domain-differentiated formal notation with a three-layer architecture that no known contemporary system matches. Whether one reads it as an exceptionally sophisticated human specialist notation or as an independently designed recording system of unknown provenance, it represents the most formally advanced notation system of its era---and the most formally interesting undeciphered text in the historical record.

% ============================================================
\section*{Acknowledgments}

The authors acknowledge the generations of researchers who have preserved, studied, and made publicly accessible the Voynich Manuscript and its transliterations. Special recognition is due to Takeshi Takahashi for his complete EVA transliteration, to Jorge Stolfi for the interlinear archive, and to Ren\'e Zandbergen for decades of scholarly curation at \url{voynich.nu}. We thank the Beinecke Rare Book and Manuscript Library at Yale University for open digital access to the manuscript.

% ============================================================
\section*{Data Availability}

All primary data, analysis scripts, and supplementary materials from this series are available at:
\url{https://github.com/RobLe3/voynich-spbceh}. The primary corpus is the Takahashi EVA transliteration, publicly available via the Voynich Manuscript Research Pages (\url{https://www.voynich.nu/}). The three papers in this series are published at \url{https://roblemumin.com/library.html}.

% ============================================================
\begin{thebibliography}{99}
\small

\bibitem{author2026a} Mumin, R. (2026a). Positional Role Asymmetry in Voynich Manuscript EVA Clusters: A Functional Classification Framework with Falsification Protocol. Available at: \url{https://roblemumin.com/library.html}.

\bibitem{author2026b} Mumin, R. (2026b). Non-Linguistic Cognitive Encoding in the Voynich Manuscript: Evidence for a State-Based Recording System from Functional Cluster Analysis. Available at: \url{https://roblemumin.com/library.html}.

\bibitem{bax2014} Bax, S. (2014). A proposed partial decoding of the Voynich script. Working paper, self-published. \url{https://stephenbax.net/?p=1175}.

\bibitem{bowern2021} Bowern, C.L., \& Lindemann, L. (2021). The linguistics of the Voynich Manuscript. \textit{Annual Review of Linguistics}, 7, 285--308. doi:10.1146/annurev-linguistics-011619-030613

\bibitem{bright1588} Bright, T. (1588). \textit{Characterie: An Arte of Short, Swifte, and Secrete Writing by Character}. London: I.\ Windet.

\bibitem{busa1980} Busa, R. (1980). The annals of humanities computing: The index Thomisticus. \textit{Computers and the Humanities}, 14, 83--90.

\bibitem{cheshire2019} Cheshire, G. (2019). The Language and Writing System of MS408 (Voynich) Explained. \textit{Romance Studies}, 37(1), 30--67.

\bibitem{currier1976} Currier, P.H. (1976). Papers on the Voynich Manuscript.

\bibitem{dimperio1978} D'Imperio, M.E. (1978). \textit{The Voynich Manuscript: An Elegant Enigma}. NSA/CSS.

\bibitem{eco1995} Eco, U. (1995). \textit{The Search for the Perfect Language}. Blackwell.

\bibitem{gonick1994} Gonick, L., \& Huffman, A. (1994). The Cartoon Guide to Statistics. Harper.

\bibitem{johannessen2000} Johannessen, J.B., et al. (2000). Codecs and annotation in language processing. \textit{Computers and the Humanities}, 34.

\bibitem{montemurro2013} Montemurro, M.A., \& Zanette, D.H. (2013). Keywords and co-occurrence patterns in the Voynich Manuscript: An information-theoretic analysis. \textit{PLoS ONE}, 8(6), e66344. doi:10.1371/journal.pone.0066344

\bibitem{reddy2011} Reddy, S., \& Knight, K. (2011). What We Know About the Voynich Manuscript. In \textit{Proceedings of the 5th ACL-HLT Workshop on Language Technology for Cultural Heritage, Social Sciences, and Humanities}, pp.\ 78--86.

\bibitem{rugg2004} Rugg, G. (2004). An Elegant Hoax? \textit{Cryptologia}, 28(1).

\bibitem{strasser2018} Strasser, B.J. (2018). \textit{Collecting Experiments: Making Big Data Biology}. University of Chicago Press.

\bibitem{willis1602} Willis, J. (1602). \textit{The Art of Stenographie}. London.

\bibitem{zandbergen2024} Zandbergen, R. (2024). Voynich MS---Research Pages. \url{https://www.voynich.nu/}.

\bibitem{zyats2016} Zyats, P., et al. (2016). Supporting Traditional Scholarship: An Investigation of the Physical Construction of the Voynich Manuscript. Yale University.

\end{thebibliography}

\end{document}
